\section{Introducción}

En un mercado turístico altamente competitivo como el de Cusco, la capacidad de respuesta rápida y la gestión eficiente de la información son pilares fundamentales para la sostenibilidad y el crecimiento de las agencias de viajes. En este contexto, Perú Flores Travel Agency Tours Operador E.I.R.L., liderada por la Sra. Silveria Flores Subileta, ha consolidado su presencia en el sector ofreciendo paquetes turísticos personalizados en destinos clave como Cusco, Arequipa y Puno.

No obstante, a pesar de su sólida base operativa y normativa, la empresa operaba bajo un modelo de madurez digital inicial, caracterizado por una marcada dependencia de procesos manuales. El flujo de trabajo diario se encontraba fragmentado en "silos de información" dispersos en herramientas no integradas como hojas de cálculo (Excel), documentos de texto, cuadernos físicos y chats de WhatsApp. Esta situación generaba tiempos de respuesta comerciales lentos, promediando entre 45 y 60 minutos por cotización, además de un riesgo operativo latente debido a la falta de una trazabilidad integral de los servicios y pasajeros.

Ante este escenario, y bajo el marco del contrato con ProInnóvate, se planteó la implementación del sistema ERP Odoo en modalidad nube (SaaS) como una solución integral. Bajo un enfoque estrictamente Out-of-the-box (funcionalidades nativas), el proyecto se centró en la digitalización de los procesos comerciales y de atención mediante los módulos de CRM, Ventas y Encuestas. Esta transformación tecnológica no solo buscó erradicar la fragmentación de datos, sino también automatizar la generación de presupuestos y establecer un sistema de retroalimentación post-viaje para fortalecer la reputación digital de la marca.

El presente informe tiene como finalidad evaluar los resultados obtenidos tras la implementación de Odoo, analizando el impacto generado en la productividad operativa, la optimización de los tiempos de atención al cliente y el nivel de adopción del sistema por parte de la organización, evidenciando los beneficios alcanzados para la escalabilidad del negocio.
